\chapter{Depersonalization}

\section*{Definition}
Depersonalization is the experience of feeling detached from one's own mind, body, or self, as if observing oneself from outside. Derealization refers to the sense that the external world is unreal, dreamlike, or distorted. These are core features of depersonalization/derealization disorder when persistent or recurrent.

\section*{Pathophysiology}
Depersonalization involves fronto-limbic dysconnectivity: reduced amygdala and insula activity with compensatory prefrontal hyperactivation, leading to emotional numbing and altered perception. The temporoparietal junction, a hub for multisensory integration and self-processing, shows abnormal connectivity, disrupting the sense of embodied self.

\section*{Causes}
\begin{itemize}
  \item \textbf{Psychiatric:} Depersonalization/derealization disorder, panic disorder, PTSD, major depressive disorder, schizophrenia
  \item \textbf{Neurological:} Temporal lobe epilepsy, migraine, traumatic brain injury, brain tumors (especially temporal lobe)
  \item \textbf{Substance-induced:} Cannabis, hallucinogens (LSD, psilocybin), ketamine, MDMA, alcohol withdrawal
  \item \textbf{Psychological:} Severe stress, trauma, sleep deprivation, meditation-induced (absorption)
\end{itemize}

\section*{When to See a Doctor}
See your doctor if:
\begin{itemize}
  \item Depersonalization or derealization persists for hours or recurs frequently
  \item Distress is significant or functioning is impaired
  \item There are associated neurological symptoms (headache, seizure-like activity, focal deficits)
\end{itemize}

\section*{Related Symptoms}
\begin{itemize}
  \item Derealization, emotional numbing, anxiety, panic attacks, distorted time perception, visual snow syndrome, déjà vu, jamais vu
\end{itemize}
\textbf{Important:} Transient depersonalization is common in the general population (up to 75\% experience mild episodes) and is often triggered by sleep deprivation or stress.

\section*{Self-Care / Home Management}
\begin{itemize}
  \item Use grounding techniques: engage the senses (holding ice, smelling strong scents, describing the environment aloud)
  \item Reduce caffeine, cannabis, and other psychoactive substances
  \item Maintain regular sleep and avoid sleep deprivation
  \item Practice mindfulness with a focus on external sensations (not body-focused)
\end{itemize}

\section*{How Your Doctor Will Evaluate You}
\begin{itemize}
  \item Clinical history differentiating depersonalization/derealization from psychosis (intact reality testing) and other dissociative disorders
  \item Structured scales: Cambridge Depersonalization Scale (CDS), Dissociative Experiences Scale (DES)
  \item Neurological evaluation (EEG, MRI) if onset is late or atypical features present
  \item Screen for substance use and medication effects
\end{itemize}

\section*{Treatment Options}
\begin{itemize}
  \item First-line: Psychoeducation and reassurance that symptoms are benign
  \item Psychotherapy: Cognitive-behavioral therapy (CBT) targeting catastrophic appraisals; mindfulness-based approaches
  \item Pharmacotherapy: No FDA-approved medications; SSRIs may help comorbid anxiety/depression; lamotrigine, naltrexone, or clonazepam have limited evidence
  \item Treat any underlying condition (e.g., epilepsy, PTSD, migraine)
\end{itemize}

\clearpage
